\chapter{Eye Strain (Asthenopia)}

\section*{Definition}
Ocular discomfort, fatigue, or pain associated with prolonged visual tasks, eye muscle imbalance, or uncorrected refractive errors.

\section*{Pathophysiology}
Prolonged near-work (digital screens, reading) causes sustained contraction of ciliary muscle (accommodation) and medial rectus muscles (convergence), leading to muscle fatigue, reduced blink rate, and ocular surface drying. Accommodative spasm may follow prolonged near-focus, causing transient distance blur.

\section*{Etiology}
\begin{itemize}
  \item Prolonged screen time (computer vision syndrome / digital eye strain)
  \item Uncorrected refractive errors (myopia, hyperopia, astigmatism, presbyopia)
  \item Poor lighting or glare
  \item Decreased blink rate while concentrating
  \item Dry eye disease
  \item Accommodative or convergence insufficiency
  \item Small font sizes or poor screen contrast
  \item Postural issues (looking up vs. down at screens)
  \item Reading in dim light
  \item Eye muscle imbalance
\end{itemize}

\section*{Indications for Evaluation}
Seek care if:
\begin{itemize}
  \item Severe or worsening symptoms
  \item Persistent eye strain despite taking breaks and correcting ergonomics, or eye strain with headaches, double vision, or vision changes
  \item Accompanied by other concerning signs
\end{itemize}

\section*{Related Symptoms}
\begin{itemize}
  \item Headache
  \item Blurred vision
  \item Dry eyes
  \item Neck pain
  \item Photophobia
\end{itemize}
\textbf{Note:} Always consider serious underlying conditions when evaluating persistent eye strain (asthenopia).

\section*{Self-Care / Home Management}
\begin{itemize}
  \item Rest and avoid activities that may aggravate the condition.
  \item Maintain adequate hydration and a balanced diet.
  \item Over-the-counter remedies may provide symptomatic relief (consult a pharmacist or doctor).
  \item Monitor symptoms and seek professional advice if they persist or worsen.
  \item Avoid self-medicating with prescription medications without medical guidance.
\end{itemize}

\section*{Diagnostic Approach}
\begin{itemize}
  \item A thorough history and physical examination are the first steps in evaluation.
  \item Laboratory tests (blood work, urinalysis) may be ordered based on clinical suspicion.
  \item Imaging studies (X-ray, ultrasound, CT, MRI) may be indicated when structural pathology is suspected.
  \item Specialist referral is arranged when the diagnosis remains uncertain or requires specific expertise.
\end{itemize}

\section*{Treatment / Management Overview}
\begin{itemize}
  \item Treatment targets the underlying cause identified through diagnostic evaluation.
  \item Supportive care and symptom management are provided while awaiting definitive therapy.
  \item Medications, lifestyle modifications, and physical therapies may be prescribed as appropriate.
  \item Follow-up ensures treatment effectiveness and allows timely adjustments to the care plan.
\end{itemize}

\clearpage